% Shared exercise chunk: l4-direct-cost-gradient
\begin{exercisechunk}
\item \textbf{Direct cost, log loss, and natural gradient.}
\label[exercise]{ex:direct-cost-gradient}
\exerciseprofile{2}{2}{2}{\exproof\quad\excalculation\quad\exinterpretation}{%
Compare three optimization dynamics and identify why an exact smooth task-cost
gradient can provide a poor training signal.}
Consider a one-step binary policy. The correct action is \(1\), and the policy
samples \(A\in\{0,1\}\) with
\[
 \pi_z(1)=p(z):=\sigma(z)=\frac{1}{1+e^{-z}},
 \qquad
 \pi_z(0)=1-p(z).
\]
The task cost is \(c(A)=\I\{A=0\}\), so its expectation is
\(C(z)=\E[c(A)]=1-p(z)\). The correct action is known, so every gradient below
is evaluated exactly. Define the log-loss surrogate
\[
 L(z)=-\log p(z).
\]
For \(a>0\), initialize \(z(0)=-a\).
For a differentiable loss \(f(z)\), its \emph{gradient flow} is the
continuous-time rule
\[
 \dot z(t)=-f'(z(t)).
\]
For this policy, define the scalar Fisher information by
\[
 F(z)=\E_{A\sim\pi_z}\!\left[
   \left(\frac{\partial}{\partial z}\log\pi_z(A)\right)^2
 \right].
\]
For a differentiable expected reward \(J(z)\), its \emph{natural-gradient
flow} for maximizing \(J\) is
\[
 \dot z_{\rm NG}(t)
 =F(z_{\rm NG}(t))^{-1}J'(z_{\rm NG}(t)).
\]
\begin{enumerate}[(a),beginpenalty=10000]
\item Use this definition to derive the gradient flows for minimizing \(C\)
and \(L\):
\[
 \dot z_C=p(z_C)(1-p(z_C)),
 \qquad
 \dot z_L=1-p(z_L).
\]
\item Let \(\tau_C(a)\) and \(\tau_L(a)\) be the first times at which the
corresponding flow reaches \(z=0\). Show that
\[
 \tau_C(a)=e^a-e^{-a}+2a,
 \qquad
 \tau_L(a)=a+1-e^{-a}.
\]
Deduce that \(\tau_C(a)\sim e^a\) and \(\tau_L(a)\sim a\) as
\(a\to\infty\).
\item Show that the Fisher information of the Bernoulli policy is
\[
 F(z)=p(z)(1-p(z)).
\]
For the expected reward \(J(z)=p(z)\), show that its natural-gradient flow is
\(\dot z_{\rm NG}(t)=1\), and hence that the flow reaches \(z=0\) at time
\(a\).
\item Explain why this example is an optimization failure rather than a
nonsmooth-loss, exploration, delayed-credit, or gradient-estimation failure.
Your explanation must also state why multiplying the direct-cost gradient by
one fixed positive constant cannot remove its exponential dependence on \(a\),
whereas log loss and natural gradient use state-dependent rescaling.
\end{enumerate}
\begin{hints}
For (b), if \(\dot z=g(z)>0\), separate variables to express the hitting time
as an integral of \(1/g(z)\).
\end{hints}
\end{exercisechunk}
